\section{Sumas sobre primos y transformadas}
\label{sec:sumas}

\fuente{A, §2, especialmente Prop.~2.1, pp.~3--4; B, §1.2,
pp.~3--5, y §2, pp.~14--19. Los cálculos de esta nota se desarrollan aquí;
el criterio completo de ventanas se demuestra en la Nota~\ref{sec:ventanas}.}

\subsection{La suma suavizada}

Sea $\chi$ un carácter primitivo de Dirichlet. La función de von Mangoldt es
\[
\Lambda(n)=
\begin{cases}
\log p,& n=p^a,\quad p\text{ primo},\quad a\ge1,\\
0,&\text{en otro caso}.
\end{cases}
\]
Para $x\ge1$ definimos
\begin{equation}
  S_\chi(x)=\sum_{n\le x}\Lambda(n)\chi(n)\left(\frac{x}{n}-1\right).
  \label{eq:suma-suave}
\end{equation}
La suma es finita. En $x=n$, el término que acaba de incorporarse vale cero:
por eso el peso suaviza los saltos de una suma de conteo. El soporte de
$\Lambda$ hace que solo contribuyan potencias de primos.

En el semiplano $\Ree s>1$, la derivación logarítmica del producto de Euler
proporciona la identidad clásica
\begin{equation}
  -\frac{L'(s,\chi)}{L(s,\chi)}
   =\sum_{n\ge1}\frac{\Lambda(n)\chi(n)}{n^s}.
  \label{eq:derivada-log}
\end{equation}
Allí las series convergen absolutamente. No se está suponiendo información
sobre ceros dentro de la franja crítica.

\subsection{Cálculo de la transformada de Mellin}

\begin{proposicion}
Para $\Ree s>1$,
\begin{equation}
  \int_1^\infty S_\chi(x)x^{-s-1}\dd x
   =-\frac{L'(s,\chi)}{s(s-1)L(s,\chi)}.
  \label{eq:mellin-s}
\end{equation}
\end{proposicion}

\begin{proof}
Escribamos $\sigma=\Ree s>1$. Para un entero $n\ge1$,
\begin{align*}
\int_n^\infty\left(\frac{x}{n}-1\right)x^{-s-1}\dd x
 &=\frac1n\int_n^\infty x^{-s}\dd x
      -\int_n^\infty x^{-s-1}\dd x\\
 &=\frac{n^{-s}}{s-1}-\frac{n^{-s}}s
  =\frac{n^{-s}}{s(s-1)}.
\end{align*}
Antes de intercambiar suma e integral comprobamos convergencia absoluta.
Como $|\chi(n)|\le1$ y $\Lambda(n)\le\log n$,
\begin{align*}
&\sum_{n\ge1}\Lambda(n)|\chi(n)|
  \int_n^\infty\left(\frac{x}{n}-1\right)x^{-\sigma-1}\dd x\\
&\hspace{2em}\le
 \frac{1}{\sigma(\sigma-1)}
 \sum_{n\ge1}\frac{\log n}{n^\sigma}<\infty.
\end{align*}
Fubini permite el intercambio. El cálculo del núcleo y
\eqref{eq:derivada-log} dan \eqref{eq:mellin-s}.
\end{proof}

La razón para elegir $x/n-1$ ya es visible: su transformada introduce el
factor explícito $1/[s(s-1)]$. Un cero de $L$ se convierte en un polo de
$L'/L$, y podremos seguirlo a través de este factor.

\subsection{La corrección en el punto \texorpdfstring{$s=1$}{s=1}}

Supongamos primero que $\chi$ no es principal. Utilizamos el resultado
clásico $L(1,\chi)\ne0$ y ponemos
\[
  \ell_\chi=\frac{L'(1,\chi)}{L(1,\chi)},\qquad
  G_\chi(x)=S_\chi(x)+x\ell_\chi.
\]
Para $\Ree s>1$, la transformada corregida es
\begin{align}
F_\chi(s)
  &:=\int_1^\infty G_\chi(x)x^{-s-1}\dd x\notag\\
  &=-\frac{L'(s,\chi)}{s(s-1)L(s,\chi)}
      +\frac{\ell_\chi}{s-1}
   =\frac{-L'(s,\chi)/L(s,\chi)+s\ell_\chi}{s(s-1)}.
   \label{eq:mellin-g}
\end{align}
El numerador se anula en $s=1$. Como $L'/L$ es holomorfa cerca de ese
punto, la singularidad es removible. \textbf{El signo positivo de
$x\ell_\chi$ es indispensable}: con el signo contrario no se produciría
esta cancelación.

Para el carácter primitivo principal, $L(s,1)=\zeta(s)$ y la corrección cambia:
\begin{equation}
 G_1(x)=S_1(x)-x\log x+(1+\gamma)x,
 \label{eq:g-principal}
\end{equation}
donde $\gamma$ es la constante de Euler. Por integración directa,
\begin{equation}
 F_1(s)=
 -\frac{\zeta'(s)}{s(s-1)\zeta(s)}
 -\frac1{(s-1)^2}+\frac{1+\gamma}{s-1}.
 \label{eq:mellin-principal}
\end{equation}
Para comprobar la cancelación, escribimos $t=s-1$ y usamos
\[
 \zeta(1+t)=\frac1t+\gamma+O(t),\qquad
 -\frac{\zeta'(1+t)}{\zeta(1+t)}=\frac1t-\gamma+O(t).
\]
El primer término de \eqref{eq:mellin-principal} es
$t^{-2}-(1+\gamma)t^{-1}+O(1)$. Los otros dos lo cancelan exactamente
hasta dejar una función holomorfa en $s=1$.

\subsection{El residuo en un cero}

\begin{proposicion}
Sea $\rho$ un cero de multiplicidad $q\ge1$ de $L(s,\chi)$, con
$0<\Ree\rho<1$. La continuación meromorfa de $F_\chi$ tiene en $\rho$
un polo simple cuyo residuo es
\begin{equation}
 \Res_{s=\rho}F_\chi(s)
 =-\frac{q}{\rho(\rho-1)}
 =\frac{q}{\rho(1-\rho)}.
 \label{eq:residuo}
\end{equation}
\end{proposicion}

\begin{proof}
Localmente, $L(s,\chi)=(s-\rho)^q u(s)$, donde $u$ es holomorfa y
$u(\rho)\ne0$. Por tanto
\[
 \frac{L'(s,\chi)}{L(s,\chi)}
 =\frac{q}{s-\rho}+\frac{u'(s)}{u(s)}.
\]
El factor $s(s-1)$ no se anula en $\rho$, y los términos de corrección
son holomorfos allí. La fórmula del residuo se sigue de
\eqref{eq:mellin-g} o \eqref{eq:mellin-principal}.
\end{proof}

Este argumento no sitúa los ceros en la recta crítica. Describe qué huella
dejaría en la transformada cualquier cero de la franja, incluso uno con
parte real mayor que $1/2$.

\subsection{Positividad y el objetivo de las ventanas}

Si $\psi$ tiene orden $o$, denotamos por $\psi^{[j]}$ el carácter primitivo
que induce $\psi^j$ y definimos
\begin{equation}
 H_\psi(x)=G_1(x)+\frac1{3o+1}
    \sum_{j=1}^{o-1}\Ree G_{\psi^{[j]}}(x).
 \label{eq:hpsi}
\end{equation}
Para $o=1$ la suma es vacía. Los caracteres retenidos forman una familia
estable bajo conjugación, de modo que la suma de sus partes reales coincide
con la suma de las funciones. En un cero compartido, todos los residuos
son múltiplos de $1/[\rho(1-\rho)]$ con multiplicidades y coeficientes
positivos. Su suma no puede cancelarse.

El objetivo aritmético puede formularse así: para cada $\eta>0$ y cada
cota inicial $A\ge2$, existe $Y\ge A$ tal que
\begin{equation}
 H_\psi(x)\le x^{1/2+\eta}
 \qquad\text{para todo }x\in[Y,Y^{3/2}].
 \label{eq:ventana}
\end{equation}
Es una cota unilateral en intervalos arbitrariamente lejanos. No se afirma
una cota de valor absoluto ni que los intervalos cubran una semirrecta.
El criterio analítico está adaptado a esos cuantificadores. La
Nota~\ref{sec:ventanas} desarrolla su demostración mediante derivadas de
la transformada de Laplace, concentración de un núcleo y selección de
un polo en discos compactos. La construcción aritmética de las ventanas
sigue pendiente de estudio. Véanse A, §2, y B, Props.~2.1--2.3.

\begin{ejercicio}
Compruebe directamente que
$\int_1^\infty x^{-s}\dd x=(s-1)^{-1}$ y
$\int_1^\infty(\log x)x^{-s}\dd x=(s-1)^{-2}$ para $\Ree s>1$.
Explique por qué estas integrales no se usan como integrales convergentes
en toda la franja crítica: allí se utiliza la continuación meromorfa de
las expresiones obtenidas.
\end{ejercicio}
